\section{Posterior updating and prediction}

\begin{qfmethod}[Bayesian updating of a discrete prior: prior, likelihood, normalising constant, posterior and predictive probability]\label{meth:bayesian-updating}
\qffield{Statement} For a partition $B_1,B_2,\dots$ of the possibilities and observed data $A$, Bayes' rule gives $P(B_j\mid A)=\frac{P(A\mid B_j)P(B_j)}{\sum_iP(A\mid B_i)P(B_i)}$, the denominator being the probability of the data by the law of total probability. A rare condition with an accurate test still leaves a low posterior probability: with prevalence 1 in 10{,}000, a 99 percent true-positive and 2 percent false-positive rate, a positive result raises the probability only to about 0.0049 \citep[sec.~1.4.3]{probability-statistics-random-processes}.

\qffield{Notation} Hypotheses $h_j$ with prior $\pi_j$; likelihood $L_j=P(\text{data}\mid h_j)$; evidence $E=\sum_j\pi_jL_j$; posterior $\pi_j'=\pi_jL_j/E$, which sums to one. Predictive probability of the next observation: $\sum_j\pi_j'P(\text{next}\mid h_j)$. Interpretation: with observations independent \emph{conditional on} the hypothesis, updating one at a time and normalising each time gives the same posterior as using all the data at once, in any order. They are not independent unconditionally: under the prior below, two right calls in a row have probability $0.25\times0.45^2+0.5\times0.5^2+0.25\times0.55^2=0.25125$, more than $0.5\times0.5=0.25$, because each success raises the probability of high-skill hypotheses.

\qffield{Conditions} The posterior is only as good as the prior and the likelihood model; here the months are assumed independent given the up-probability, which is constant, and the true value is assumed to be one of the listed hypotheses. Interpretation: a prior concentrated on few values cannot learn values it excludes.

\qffield{Practice} Example, original and synthetic, not evidence about any strategy. The monthly probability that a signal is right is 0.45, 0.50 or 0.55, with prior probabilities 0.25, 0.50 and 0.25. In 20 independent months it is right 14 times. The binomial likelihoods $\binom{20}{14}h^{14}(1-h)^6$ are 0.01498, 0.03696 and 0.07460; weighted by the priors they sum to the evidence 0.04088; the posterior is 0.0916, 0.4521 and 0.4562. The predictive probability of a right call next month rises from 0.5000 to 0.5182, well below the sample frequency of 0.70, because the prior allows nothing above 0.55. Updating month by month gives the same posterior.

\qffield{Pitfalls} Omitting the denominator (the posterior then does not sum to one); confusing $P(\text{data}\mid h)$ with $P(h\mid\text{data})$; treating the sample frequency as the posterior. Choosing the prior after seeing the data.

\qffield{See also} \hyperref[def:conditional-moments]{conditional moments}; \hyperref[meth:skill-luck-significance-tests]{skill against luck} (a frequentist test of the same question).

\end{qfmethod}
